\documentclass[12pt,CJK]{article}
\usepackage{geometry}
\usepackage{ulem}
\usepackage{CJKutf8}
\usepackage{amsmath}
\usepackage{amssymb}
\usepackage{graphicx}
\usepackage{eufrak}
\usepackage{color}
\usepackage{slashed}
\usepackage{simplewick}
\usepackage{tikz}
\usepackage{tcolorbox}
\graphicspath{{../figures/}}
%%figures
\def\lfig#1#2{\includegraphics[width=#1 in]{#2}}
\def\addfig#1#2{\begin{center}\includegraphics[width=#1 in]{#2}\end{center}}
\def\wulian{\includegraphics[width=0.18in]{emoji_wulian.jpg}}
\def\bigwulian{\includegraphics[width=0.35in]{emoji_wulian.jpg}}
\def\bye{\includegraphics[width=0.18in]{emoji_bye.jpg}}
\def\bigbye{\includegraphics[width=0.35in]{emoji_bye.jpg}}
\def\huaixiao{\includegraphics[width=0.18in]{emoji_huaixiao.jpg}}
\def\bighuaixiao{\includegraphics[width=0.35in]{emoji_huaixiao.jpg}}
\def\jianxiao{\includegraphics[width=0.18in]{emoji_jianxiao.jpg}}
\def\bigjianxiao{\includegraphics[width=0.35in]{emoji_jianxiao.jpg}}
%% colors
\def\blacktext#1{{\color{black}#1}}
\def\bluetext#1{{\color{blue}#1}}
\def\redtext#1{{\color{red}#1}}
\def\darkbluetext#1{{\color[rgb]{0,0.2,0.6}#1}}
\def\skybluetext#1{{\color[rgb]{0.2,0.7,1.}#1}}
\def\cyantext#1{{\color[rgb]{0.,0.5,0.5}#1}}
\def\greentext#1{{\color[rgb]{0,0.7,0.1}#1}}
\def\darkgray{\color[rgb]{0.2,0.2,0.2}}
\def\lightgray{\color[rgb]{0.6,0.6,0.6}}
\def\gray{\color[rgb]{0.4,0.4,0.4}}
\def\blue{\color{blue}}
\def\red{\color{red}}
\def\green{\color{green}}
\def\darkgreen{\color[rgb]{0,0.4,0.1}}
\def\darkblue{\color[rgb]{0,0.2,0.6}}
\def\skyblue{\color[rgb]{0.2,0.7,1.}}
%%control
\def\be{\begin{equation}}
\def\ee{\nonumber\end{equation}}
\def\bea{\begin{eqnarray}}
\def\eea{\nonumber\end{eqnarray}}
\def\bch{\begin{CJK}{UTF8}{gbsn}}
\def\ech{\end{CJK}}
\def\bitem{\begin{itemize}}
\def\eitem{\end{itemize}}
\def\bcenter{\begin{center}}
\def\ecenter{\end{center}}
\def\bex{\begin{minipage}{0.2\textwidth}\includegraphics[width=0.6in]{jugelizi.png}\end{minipage}\begin{minipage}{0.76\textwidth}}
\def\eex{\end{minipage}}
\def\chtitle#1{\frametitle{\bch#1\ech}}
\def\bmat#1{\left(\begin{array}{#1}}
\def\emat{\end{array}\right)}
\def\bcase#1{\left\{\begin{array}{#1}}
\def\ecase{\end{array}\right.}
\def\bmini#1{\begin{minipage}{#1\textwidth}}
\def\emini{\end{minipage}}
\def\tbox#1{\begin{tcolorbox}#1\end{tcolorbox}}
\def\pfrac#1#2#3{\left(\frac{\partial #1}{\partial #2}\right)_{#3}}
%%symbols
\def\bropt{\,(\ \ \ )}
\def\optlist#1#2#3{(A)\;#1\ \ \ (B)\;#2\ \ \ (C)\;#3}
\def\sone{$\star$}
\def\stwo{$\star\star$}
\def\sthree{$\star\star\star$}
\def\sfour{$\star\star\star\star$}
\def\sfive{$\star\star\star\star\star$}
\def\rint{{\int_\leftrightarrow}}
\def\roint{{\oint_\leftrightarrow}}
\def\stdHf{{\textit{\r H}_f}}
\def\deltaH{{\Delta \textit{\r H}}}
\def\ii{{\dot{\imath}}}
\def\skipline{{\vskip0.1in}}
\def\skiplines{{\vskip0.2in}}
\def\lagr{{\mathcal{L}}}
\def\hamil{{\mathcal{H}}}
\def\vecv{{\mathbf{v}}}
\def\vecx{{\mathbf{x}}}
\def\vecy{{\mathbf{y}}}
\def\veck{{\mathbf{k}}}
\def\vecp{{\mathbf{p}}}
\def\vecn{{\mathbf{n}}}
\def\vecA{{\mathbf{A}}}
\def\vecP{{\mathbf{P}}}
\def\vecsigma{{\mathbf{\sigma}}}
\def\hatJn{{\hat{J_\vecn}}}
\def\hatJx{{\hat{J_x}}}
\def\hatJy{{\hat{J_y}}}
\def\hatJz{{\hat{J_z}}}
\def\hatj#1{\hat{J_{#1}}}
\def\hatphi{{\hat{\phi}}}
\def\hatq{{\hat{q}}}
\def\hatpi{{\hat{\pi}}}
\def\vel{\upsilon}
\def\Dint{{\mathcal{D}}}
\def\adag{{\hat{a}^\dagger}}
\def\bdag{{\hat{b}^\dagger}}
\def\cdag{{\hat{c}^\dagger}}
\def\ddag{{\hat{d}^\dagger}}
\def\hata{{\hat{a}}}
\def\hatb{{\hat{b}}}
\def\hatc{{\hat{c}}}
\def\hatd{{\hat{d}}}
\def\hatN{{\hat{N}}}
\def\hatH{{\hat{H}}}
\def\hatp{{\hat{p}}}
\def\Fup{{F^{\mu\nu}}}
\def\Fdown{{F_{\mu\nu}}}
\def\newl{\nonumber \\}
\def\vece{\mathrm{e}}
\def\calM{{\mathcal{M}}}
\def\calT{{\mathcal{T}}}
\def\calR{{\mathcal{R}}}
\def\barpsi{\bar{\psi}}
\def\baru{\bar{u}}
\def\barv{\bar{\upsilon}}
\def\qeq{\stackrel{?}{=}}
\def\torder#1{\mathcal{T}\left(#1\right)}
\def\rorder#1{\mathcal{R}\left(#1\right)}
\def\contr#1#2{\contraction{}{#1}{}{#2}#1#2}
\def\trof#1{\mathrm{Tr}\left(#1\right)}
\def\trace{\mathrm{Tr}}
\def\comm#1{\ \ \ \left(\mathrm{used}\ #1\right)}
\def\tcomm#1{\ \ \ (\text{#1})}
\def\slp{\slashed{p}}
\def\slk{\slashed{k}}
\def\calp{{\mathfrak{p}}}
\def\veccalp{\mathbf{\mathfrak{p}}}
\def\Tthree{T_{\tiny \textcircled{3}}}
\def\pthree{p_{\tiny \textcircled{3}}}
\def\dbar{{\,\mathchar'26\mkern-12mu d}}
\def\erf{\mathrm{erf}}
\def\const{\mathrm{constant}}
\def\pheat{\pfrac p{\ln T}V}
\def\vheat{\pfrac V{\ln T}p}
%%units
\def\fdeg{{^\circ \mathrm{F}}}
\def\cdeg{^\circ \mathrm{C}}
\def\atm{\,\mathrm{atm}}
\def\angstrom{\,\text{\AA}}
\def\SIL{\,\mathrm{L}}
\def\SIkm{\,\mathrm{km}}
\def\SIyr{\,\mathrm{yr}}
\def\SIGyr{\,\mathrm{Gyr}}
\def\SIV{\,\mathrm{V}}
\def\SImV{\,\mathrm{mV}}
\def\SIeV{\,\mathrm{eV}}
\def\SIkeV{\,\mathrm{keV}}
\def\SIMeV{\,\mathrm{MeV}}
\def\SIGeV{\,\mathrm{GeV}}
\def\SIcal{\,\mathrm{cal}}
\def\SIkcal{\,\mathrm{kcal}}
\def\SImol{\,\mathrm{mol}}
\def\SIN{\,\mathrm{N}}
\def\SIHz{\,\mathrm{Hz}}
\def\SIm{\,\mathrm{m}}
\def\SIcm{\,\mathrm{cm}}
\def\SIfm{\,\mathrm{fm}}
\def\SImm{\,\mathrm{mm}}
\def\SInm{\,\mathrm{nm}}
\def\SImum{\,\mathrm{\mu m}}
\def\SIJ{\,\mathrm{J}}
\def\SIW{\,\mathrm{W}}
\def\SIkJ{\,\mathrm{kJ}}
\def\SIs{\,\mathrm{s}}
\def\SIkg{\,\mathrm{kg}}
\def\SIg{\,\mathrm{g}}
\def\SIK{\,\mathrm{K}}
\def\SImmHg{\,\mathrm{mmHg}}
\def\SIPa{\,\mathrm{Pa}}

\geometry{tmargin=0.5in, bmargin=0.5in, lmargin=0.6in, rmargin=0.9in, nohead, nofoot}
\def\mark#1{{\ \color{blue}(#1分)}}
\renewcommand{\thepage}{}
\begin{document}
\bch
\begin{center}
数学物理方法2019 课堂小测II

{\vskip 0.1in}

  共100分： 30分签到 + 70分 piece of cake 

\end{center}
{\vskip 0.1in}

{\bf 请在答题纸上任意选做下列题目（标明题号），\sout{凑够70分就行}\ 能凑几分算几分吧……}

{\vskip 0.1in}

\bitem
\item[(1)]{(为了检验你期中以前内容是否都量子波动速读了）请计算沿着逆时针方向的围道的积分：
    $$\int_{|z|=1}\, \frac{1}{\sin z - z\cos z} dz.$$
  注意使用留数定理时不要漏掉孤立奇点。\mark{15}}
\item[(2)]{设函数 $f(x,y), g(x,y)$ 都定义在直角坐标系的椭圆 $\Omega = \{(x,y): \frac{x^2}{4}+y^2\le 1\} $ 内，且处处满足方程
    $$  (x^2y^2 + \nabla^2)f(x,y) = (1-x^2y^2-\nabla^2)g(x,y)  = 0.$$
    在$\Omega$的边界上, $f$ 和 $g$ 满足：
    $$\left.x\frac{\partial f}{\partial x}+4y\frac{\partial f}{\partial y}\right\vert_{\frac{x^2}{4}+y^2=1} =  \left.x\frac{\partial g}{\partial x}+4y\frac{\partial g}{\partial y}\right\vert_{\frac{x^2}{4}+y^2=1} = 0.$$
    证明： $$\int_\Omega f(x,y)g(x,y) dx dy = 0.$$ \mark{20}
  }  
\item[(3)]{设函数 $f(x) = x^\alpha J_m(x)$ （这里的$\alpha$为实数, $m$为非负整数）在 $x=\mu$ 和 $x=\nu$ 均取到极值，且 $\mu>\nu>0$。请证明 $\int_0^1 x J_m(\mu x)J_m(\nu x) dx  = 0.$ \mark{15}}
\item[(4)]{两端固定在同一水平线上，线质量密度为$\rho$，长度为 $L$ 的绷直的均匀弦上的张力为$F$。把弦的中点竖直往上拉$\epsilon$的小距离，然后从静止释放。忽略重力、空气阻力的影响，弦之后如何振动？ \mark{25}}
\item[(5)]{一个半径为 $R$ 的均匀孤立圆盘，其材质的热传导方程参数为 $a$。一开始圆盘的温度分布为 $$\left.T\right\vert_{t=0} = T_0\left(2-\frac{r^2}{R^2}\right)\cos^2\theta,$$这里的 $(r,\theta)$ ($0\le r\le R$) 是以圆盘中心为原点建立的极坐标，$T_0$ 为常量。请计算之后圆盘上的温度变化 $T(r,\theta, t)$。\mark{30}}
\item[(6)]{一个半径为 $R$，周边固定的均匀圆形薄膜，其横向小振动的波速参数为 $a$。一开始薄膜各处均处于静止，横向小位移为 $$\left.u\right\vert_{t=0} = A\frac{r}{R}\left(1-\frac{r^2}{R^2}\right)\cos\theta,$$这里的 $(r,\theta)$ ($0\le r\le R$) 是以薄膜中心为原点建立的极坐标，$A$为常量。请计算之后薄膜的横向小振动 $u(r,\theta, t)$。\mark{30}}
\item[(7)]{从前有一个神秘的星球叫浮士德星，这个星球的每年有一万天。星球上的人出生时都会和魔鬼签署了一份协议：每个人每天起床都要投掷 $600$ 次骰子。如果 $1,2,3,4,5,6$ 每个数字均恰好出现 $100$ 次，魔鬼就要带走这个人的灵魂，否则魔鬼会保护这个人一天平安无事。\\
    有一天赌王来到了浮士德星，传授给了浮士德星人一个秘术——666大法。使用秘术之后，每次掷骰子出现 $6$ 的概率会稍稍增加，但增加的幅度非常小以致于魔鬼无法察觉。从此浮士德星人的平均寿命延长了5年。请由此推断: 秘术能使单次掷骰子出现 $6$ 的概率增加多少？ （结果请保留至少2位有效数） \mark{30}}
\item[(8)]{一个质量为$m$，固有圆频率为 $\omega$ 的谐振子，其三个空间坐标 $x,y,z$ 和三个动量分量 $p_x, p_y, p_z$ 构成了六维的 ``相空间''。它在任意时刻的状态（即位置和运动情况）可以用相空间的一个点来表示。从经典物理的角度看，谐振子的状态允许是相空间中任意一个点。但如果要求谐振子的能量不超过某个阈值 $E_0$ (常量$E_0>0$)，即
    $$\frac{1}{2}m\omega^2 \left(x^2+y^2+z^2\right) + \frac{1}{2m}\left(p_x^2+p_y^2+p_z^2\right) \le E_0,$$      
    谐振子就被限制在相空间一个有限大小的区域内。这个区域的（六维）体积是多少？\mark{15}}
\item[(9)]{假设你MMP考试挂科，未能毕业，流落江湖摆摊卖糖葫芦为生。六一儿童节到了，糖葫芦摊抽奖送礼大酬宾。你在一个大袋子里装了很多红、黄、蓝色的小球，每次摸出每种颜色小球的概率都是1/3。你请光临摊位的小孩随机摸出 12 个小球，不计颜色的次序，摸出的三种颜色小球的个数有18种可能性: (0,0,12), (0,1,11), (0,2,10), (0,3,9), (0,4,8), (0,5,7), (0,6,6), (1,1,10), (1,2,9), (1,3,8), (1,4,7), (1,5,6), (2,2,8), (2,3,7), (2,4,6), (2,5,5), (3,4,5), (4,4,4)。只要小孩摸出的不是 (3,4,5)，就白送一根糖葫芦；否则需要按日常价 10 元购买一根糖葫芦。问：为了在大酬宾活动中不亏本，平均一根糖葫芦的成本价不能高于多少元？（精确到0.01元) \mark{10}}
\item[(10)]{已知 $x$ 为正实数，且 $\sum_{n=0}^\infty J_{2n}(x) = 0.98$，估算 $x$ 的值（至少精确到2位有效数）。 \mark{20}}
\item[(11)]{传古代印度与中国之间有一萨桑国，国王山鲁亚尔生性残暴嫉妒，因王后行为不端，将其杀死，此后每日娶一少女，翌日晨即杀掉，以示报复。宰相的女儿山鲁佐德为拯救无辜的女子，自愿嫁给国王，用让国王刷quiz的方法吸引国王，每夜刷到最精彩处，天刚好亮了，使国王爱不忍杀。quiz一直刷了一千零一夜，国王终于被感动，与她白首偕老。 \\
山鲁佐德给国王传授了自然对数的算法。国王深深迷醉于这个超越时代近千年的技术。每天刷quiz前，国王都要练习自然对数的计算。\\
为了让国王每天有新的挑战，山鲁佐德使用了满足下列递推式的数列：
$$a_n = \frac{4}{3+a_{n-1}^2}.$$
数列的起始点设为$a_0=0$，也就是第一天给国王的数字是$a_1=\frac{4}{3}$，之后依次递推。\\
国王在第$n$天（$n=1,2,...,1001$)刷quiz前的热身运动就是计算$a_n$的对数的绝对值的对数，也就是$\ln|\ln a_n|$。在当时可没有计算机，全靠手算，山鲁佐德只要求国王算到一位有效数就可以了。\\
那么问题来了：在最后第一千零一夜，国王算出来的$\ln|\ln a_{1001}|$等于多少？ \mark{20} }
\item[(12)]{计算 $J_2(x)+ xJ_2'(x)=0$ 的所有正实数解的四次方的倒数之和。 \mark{20} }
\item[(13)]{有一个半径为 $R=0.1\mathrm{m}$，材质的热传导方程参数为 $a=10^{-4}\mathrm{m^2s^{-1}}$ 的孤立、均匀薄圆盘。初始时刻上半个圆盘 (极坐标 $0\le \theta < \pi$) 的温度为 $400\mathrm{K}$，下半个圆盘 （极坐标 $\pi \le\theta < 2\pi$）的温度为 $300\mathrm{K}$。经过一分钟之后，上半圆盘的最高边缘，即 $r=R, \theta=\frac{\pi}{2}$ 处，温度是多少？(请计算至少精确到$1\mathrm{K}$) \mark{70}}  
  \eitem

  
  
\ech
\end{document}
